% Use xelatex to generate a pdf file of this presentation

\documentclass{beamer}
\usepackage{roboto}
\usepackage[english]{babel}
\usetheme{Copenhagen}

% Simple way to put a code listing to the slide
\usepackage{listings}

\usepackage{tikz}
\usepackage{graphicx}

% Blocks settings
\setbeamerfont{block title}{size=\tiny}

%Information to be included in the title page:
\title{Enforcing Pre-sorted Scans}
\subtitle{on Query Pathkeys in PostgreSQL 18}
\author{Andrei Lepikhov}
\institute{pgEdge}
\date{pgconf.dev 2026}

% Add page numbers in the footnote
\expandafter\def\expandafter\insertshorttitle\expandafter{%
  \insertshorttitle\hfill%
  \insertframenumber\,/\,\inserttotalframenumber}

\definecolor{lightlightgray}{gray}{0.9}

% Default listing format
\lstset{language=sql, frame=none, tabsize=2, identifierstyle=\color{black},
  showspaces=false, showtabs=false, showstringspaces=false,
  columns=fullflexible, % Restore default spaces between letters
  basicstyle=\rmfamily\scriptsize,
  stringstyle=\color{purple},
  keywordstyle=\bfseries\color{green!40!black},
  numberstyle=\tiny\color{gray}}

% -----------------------------

\begin{document}

% %%%%%%%%%%%%%%%%%%%%%%%%%%%%%%%%%%%%%%%%%%%%%%%%%%%%%%%%%%%%%%%%%%%%%%%%
% Enforcing Pre-sorted Scans on Query Pathkeys
% pgconf.dev 2026
% %%%%%%%%%%%%%%%%%%%%%%%%%%%%%%%%%%%%%%%%%%%%%%%%%%%%%%%%%%%%%%%%%%%%%%%%

\begin{frame}
\titlepage
\end{frame}

% %%%%%%%%%%%%%%%%%%%%%%%%%%%%%%%%%%%%%%%%%%%%%%%%%%%%%%%%%%%%%%%%%%%%%%%%
% Self Introduction
% %%%%%%%%%%%%%%%%%%%%%%%%%%%%%%%%%%%%%%%%%%%%%%%%%%%%%%%%%%%%%%%%%%%%%%%%

\begin{frame}[fragile]\frametitle{Self Introduction}
\begin{itemize}
  \item 2007 - First touch to PostgreSQL (distributed query execution)
  \item 2010 - designer and user of databases based on PostgreSQL (and SQLite)
  \item 2017 - now - PostgreSQL Enthusiast, sporadic contributor and patch reviewer
\end{itemize}
My projects:
\begin{itemize}
  \item Multimaster, Shardman, Joinsel
  \item AQO, sr\_plan, SwitchJoin, re-optimisation ...
  \item Patches: GROUP-BY, Self-Join Removal, OR->ANY Optimisation ...
\end{itemize}
\end{frame}

% %%%%%%%%%%%%%%%%%%%%%%%%%%%%%%%%%%%%%%%%%%%%%%%%%%%%%%%%%%%%%%%%%%%%%%%%
% The Problem
% %%%%%%%%%%%%%%%%%%%%%%%%%%%%%%%%%%%%%%%%%%%%%%%%%%%%%%%%%%%%%%%%%%%%%%%%

\begin{frame}[fragile]\frametitle{The Problem -- A Real-Life Scenario}
% A typical sales system query: show top-10 most popular products
% with type-specific properties from multiple tables.
\begin{lstlisting}
SELECT p.id, p.name, p.category, p.popularity,
  json_strip_nulls(json_build_object(
    'warranty', e.warranty_months,
    'voltage',  e.voltage,
    'size',     c.size,
    'color',    c.color,
    'expiry',   f.expiry_days,
    'organic',  f.organic
  )) AS properties
FROM products p                                -- 100k rows
  LEFT JOIN electronics_props e ON e.product_id = p.id
  LEFT JOIN clothing_props    c ON c.product_id = p.id
  LEFT JOIN food_props        f ON f.product_id = p.id
ORDER BY p.popularity DESC
LIMIT 10;
\end{lstlisting}
\medskip
\textbf{Key:} \texttt{products(id)} has a PK index (needed for joins), but
\textit{no index on} \texttt{popularity} --- it changes with every sale,
maintaining such an index is impractical.
\end{frame}

\begin{frame}[fragile]\frametitle{The Problem -- PG17 Plan}
% The planner joins everything first, then sorts the entire result.
\begin{lstlisting}
 Limit
   ->  Sort
         Sort Key: p.popularity DESC
         ->  Nested Loop Left Join
               ->  Nested Loop Left Join
                     ->  Nested Loop Left Join
                           ->  Seq Scan on products p
                           ->  Index Scan ... on electronics_props e
                     ->  Index Scan ... on clothing_props c
               ->  Index Scan ... on food_props f
\end{lstlisting}
\bigskip
No index on the sort column $\Rightarrow$ full sort of \textbf{100k joined rows}
just to return \textbf{10}.
\end{frame}

% %%%%%%%%%%%%%%%%%%%%%%%%%%%%%%%%%%%%%%%%%%%%%%%%%%%%%%%%%%%%%%%%%%%%%%%%
% Background: Sorting in the Planner Today
% %%%%%%%%%%%%%%%%%%%%%%%%%%%%%%%%%%%%%%%%%%%%%%%%%%%%%%%%%%%%%%%%%%%%%%%%

\begin{frame}\frametitle{Background: Pathkeys and Sort Order}
% Brief refresher:
% - Pathkeys represent sort order of a path
% - Equivalence classes group expressions that are equal
% - The planner generates sorted paths via index scans, explicit Sort, incremental sort
\end{frame}

\begin{frame}\frametitle{Background: The Gap}
% The planner never generates Sort + SeqScan for a base relation
% just because it might help a join higher in the plan tree.
% No "lookahead" from base-relation planning to top-level ORDER BY.
\end{frame}

% %%%%%%%%%%%%%%%%%%%%%%%%%%%%%%%%%%%%%%%%%%%%%%%%%%%%%%%%%%%%%%%%%%%%%%%%
% The Idea
% %%%%%%%%%%%%%%%%%%%%%%%%%%%%%%%%%%%%%%%%%%%%%%%%%%%%%%%%%%%%%%%%%%%%%%%%

\begin{frame}\frametitle{The Idea: Push Sort Below the Join}
% If a prefix of query_pathkeys can be computed from a single relation,
% generate a pre-sorted scan path for that relation.
% The NestLoop preserves sort order -> top-level Sort disappears.
% With LIMIT, the Sort becomes bounded top-N heapsort.
\end{frame}

\begin{frame}[fragile]\frametitle{The Idea -- Before and After}
\begin{columns}[T]\begin{column}{0.5\textwidth}
\begin{block}{PG17: Sort on top}
\begin{lstlisting}[basicstyle=\tiny]
 Limit
   -> Sort  <-- sorts 100k rows!
     Sort Key: p.popularity DESC
     -> NL Left Join
       -> NL Left Join
         -> NL Left Join
           -> Seq Scan on products p
           -> IdxScan electronics_props
         -> IdxScan clothing_props
       -> IdxScan food_props
\end{lstlisting}
\end{block}
\end{column}\begin{column}{0.5\textwidth}
\begin{block}{PG18: Sort pushed down}
\begin{lstlisting}[basicstyle=\tiny]
 Limit
   -> NL Left Join
     -> NL Left Join
       -> NL Left Join
         -> Sort  <-- top-N heapsort!
           Sort Key: p.popularity DESC
           -> Seq Scan on products p
         -> IdxScan electronics_props
       -> IdxScan clothing_props
     -> IdxScan food_props
\end{lstlisting}
\end{block}
\end{column}\end{columns}
\bigskip
\begin{center}
No index on the sort column --- but with LIMIT 10 the Sort becomes
a \textbf{bounded top-N heapsort} over the base table only.
\end{center}
\end{frame}

% %%%%%%%%%%%%%%%%%%%%%%%%%%%%%%%%%%%%%%%%%%%%%%%%%%%%%%%%%%%%%%%%%%%%%%%%
% Implementation
% %%%%%%%%%%%%%%%%%%%%%%%%%%%%%%%%%%%%%%%%%%%%%%%%%%%%%%%%%%%%%%%%%%%%%%%%

\begin{frame}
\vspace*{\fill}\begin{center}
\huge{Implementation}
\end{center}\vspace*{\fill}
\end{frame}

\begin{frame}[fragile]\frametitle{generate\_presorted\_paths()}
% Called from set_plain_rel_pathlist() for each base relation.
% Algorithm:
% 1. Check enable_presorted_scan GUC and query_pathkeys existence
% 2. Walk query_pathkeys prefix, check relation_can_be_sorted_early()
% 3. If useful pathkeys found and no path with them exists yet:
%    - copy_path_for_sort() the cheapest unsorted path
%    - Wrap with create_sort_path()
%    - add_path() to the relation's pathlist
\end{frame}

\begin{frame}[fragile]\frametitle{relation\_can\_be\_sorted\_early()}
% Safety gate -- same as used for FDW/custom-scan sorted paths.
% Rejects:
% - Volatile expressions
% - Set-returning functions
% - Expressions not computable from the relation's reltarget
% - Non-parallel-safe expressions
\end{frame}

\begin{frame}[fragile]\frametitle{copy\_path\_for\_sort()}
% Creates an independent shallow copy of the path.
% Same technique as reparameterize_path().
% Supports: SeqScan, SampleScan, IndexPath, BitmapHeapPath, TidPath, CustomPath.
% Prevents the Sort subpath from sharing state with the original in pathlist.
\end{frame}

\begin{frame}[fragile]\frametitle{Tuple Bound Propagation Through NestLoop}
% ExecSetTupleBound() changes:
% LEFT/ANTI joins -> propagate bound to OUTER side
% RIGHT joins -> propagate bound to INNER side
% These join types guarantee >= 1 output row per driving-side row,
% so a LIMIT-derived bound safely applies -> enables top-N heapsort.
%
% This is the EXECUTOR part of the feature -- it lives in ExecSetTupleBound(),
% which no extension hook can reach.
\end{frame}

\begin{frame}[fragile]\frametitle{The enable\_presorted\_scan GUC}
% New boolean GUC, default true.
% PGC_USERSET -- per-session control.
% GUC_EXPLAIN flag -- visible in EXPLAIN output.
% SET enable_presorted_scan = off; -- to disable
\end{frame}

% %%%%%%%%%%%%%%%%%%%%%%%%%%%%%%%%%%%%%%%%%%%%%%%%%%%%%%%%%%%%%%%%%%%%%%%%
% Challenges and Edge Cases
% %%%%%%%%%%%%%%%%%%%%%%%%%%%%%%%%%%%%%%%%%%%%%%%%%%%%%%%%%%%%%%%%%%%%%%%%

\begin{frame}
\vspace*{\fill}\begin{center}
\huge{Challenges \& Edge Cases}
\end{center}\vspace*{\fill}
\end{frame}

\begin{frame}\frametitle{Equivalence Classes Spanning Joins}
% Problem: EC for t1.x = t2.x has ec_relids = {t1, t2}.
% Naive bms_is_subset(ec->ec_relids, rel->relids) rejects it,
% even though the EC contains a member computable from t1 alone.
% Solution: looser pre-filter + explicit ec_members scan.
\end{frame}

\begin{frame}\frametitle{Volatile and Unsafe Expressions}
% Volatile ECs, SRFs, partial functions (e.g., regclass casts)
% must not be sorted early.
% relation_can_be_sorted_early() handles all of these.
\end{frame}

\begin{frame}\frametitle{Adding Sort Expressions to Pathtarget}
% Initially refused to add expressions not already in the scan's pathtarget.
% Too conservative -- disabled the optimization for most non-trivial ORDER BY.
% SQL standard: expression evaluation order is implementation-dependent.
% Core planner already does the same via resjunk injection.
\end{frame}

\begin{frame}\frametitle{Why Core, Not an Extension?}
% The planner part (generate_presorted_paths) could in theory use
% the set_rel_pathlist_hook or custom path machinery.
% But the feature has TWO parts:
%
% 1. Planner: generate pre-sorted paths (hookable, in principle)
% 2. Executor: propagate tuple bound through NestLoop in ExecSetTupleBound()
%
% Without the executor change, LIMIT cannot be pushed down to the Sort node
% below the join -> no bounded top-N heapsort -> much of the benefit is lost.
%
% ExecSetTupleBound() has no extension hook.
% This makes the feature impossible to maintain as a standalone extension.
\end{frame}

% %%%%%%%%%%%%%%%%%%%%%%%%%%%%%%%%%%%%%%%%%%%%%%%%%%%%%%%%%%%%%%%%%%%%%%%%
% Results
% %%%%%%%%%%%%%%%%%%%%%%%%%%%%%%%%%%%%%%%%%%%%%%%%%%%%%%%%%%%%%%%%%%%%%%%%

\begin{frame}
\vspace*{\fill}\begin{center}
\huge{Results}
\end{center}\vspace*{\fill}
\end{frame}

\begin{frame}[fragile]\frametitle{Products Example -- enable\_presorted\_scan = on}
% Actual EXPLAIN ANALYZE output from PG18 branch.
\begin{lstlisting}[basicstyle=\tiny]
 Limit (actual time=43.924..44.176 rows=10 loops=1)
   Buffers: shared hit=940
   ->  Nested Loop Left Join (actual time=43.922..44.172 rows=10 loops=1)
         ->  Nested Loop Left Join (actual time=43.865..43.980 rows=10 loops=1)
               ->  Nested Loop Left Join (actual time=43.851..43.903 rows=10 loops=1)
                     ->  Sort (actual time=43.825..43.826 rows=10 loops=1)
                           Sort Key: p.popularity DESC
                           Sort Method: top-N heapsort  Memory: 26kB
                           ->  Seq Scan on products p (actual rows=100000 loops=1)
                     ->  Index Scan ... on electronics_props e
               ->  Index Scan ... on clothing_props c
         ->  Index Scan ... on food_props f
\end{lstlisting}
\begin{center}
\textbf{44 ms} --- Sort is \textit{below} the joins, top-N heapsort with LIMIT 10.
\end{center}
\end{frame}

\begin{frame}[fragile]\frametitle{Products Example -- enable\_presorted\_scan = off}
% Same query with the optimization disabled.
\begin{lstlisting}[basicstyle=\tiny]
 Limit (actual time=813.491..813.495 rows=10 loops=1)
   Buffers: shared hit=1442
   ->  Sort (actual time=813.489..813.491 rows=10 loops=1)
         Sort Key: p.popularity DESC
         Sort Method: top-N heapsort  Memory: 27kB
         ->  Hash Left Join (actual time=73.483..772.134 rows=100000 loops=1)
               ->  Hash Left Join (actual time=54.732..128.615 rows=100000 loops=1)
                     ->  Hash Left Join (actual time=32.330..79.951 rows=100000 loops=1)
                           ->  Seq Scan on products p (actual rows=100000 loops=1)
                           ->  Hash (rows=33333)
                                 ->  Seq Scan on electronics_props e
                     ->  Hash (rows=33334)
                           ->  Seq Scan on clothing_props c
               ->  Hash (rows=33333)
                     ->  Seq Scan on food_props f
\end{lstlisting}
\begin{center}
\textbf{814 ms} --- Hash joins all 100k rows, \textit{then} sorts. \textbf{18$\times$ slower.}
\end{center}
\end{frame}

\begin{frame}\frametitle{Regression Test Impact}
% 8+ regression tests changed plans:
% - join, incremental_sort, groupingsets, collate.icu.utf8,
%   partition_aggregate, partition_join, partition_prune, tuplesort
% + contrib: pg_overexplain, postgres_fdw
%
% Demonstrates the optimizer picks up the optimization in diverse real patterns.
\end{frame}

% %%%%%%%%%%%%%%%%%%%%%%%%%%%%%%%%%%%%%%%%%%%%%%%%%%%%%%%%%%%%%%%%%%%%%%%%
% Future Work
% %%%%%%%%%%%%%%%%%%%%%%%%%%%%%%%%%%%%%%%%%%%%%%%%%%%%%%%%%%%%%%%%%%%%%%%%

\begin{frame}\frametitle{Future Work}
% - Extending to other path types (ForeignScan, Append)
% - Interaction with parallel queries
% - Cost model refinements
% - Benchmarks on real-world workloads (1C:Enterprise, TPC-H)
% - Community feedback and review status
\end{frame}

% %%%%%%%%%%%%%%%%%%%%%%%%%%%%%%%%%%%%%%%%%%%%%%%%%%%%%%%%%%%%%%%%%%%%%%%%
% Closing
% %%%%%%%%%%%%%%%%%%%%%%%%%%%%%%%%%%%%%%%%%%%%%%%%%%%%%%%%%%%%%%%%%%%%%%%%

\begin{frame}
\vspace*{\fill}
\begin{center}
\huge{Questions ?}
\end{center}
\vspace*{\fill}
\end{frame}

\end{document}
